\chapter*{Introduction générale}
\addcontentsline{toc}{chapter}{Introduction générale}
\markboth{Introduction générale}{Introduction générale}

Les dispositifs médicaux respiratoires, tels que les concentrateurs d'oxygène, les appareils de ventilation non invasive (VNI) et les appareils à pression positive continue (CPAP), sont utilisés par des patients dont la santé dépend directement de leur bon fonctionnement. Leur maintenance doit donc être rigoureuse, mais aussi \textbf{traçable} : savoir à tout moment où se trouve un appareil, qui est intervenu, quelle panne a été constatée, quelles pièces ont été remplacées et si l'appareil a été contrôlé avant d'être remis à disposition.

Au sein du service Maintenance d'Air Liquide Tunisie, cette traçabilité reposait, au début du stage, sur des documents papier remplis pendant la journée, puis sur un classeur Google Sheets alimenté en fin de journée. Cette organisation, simple à mettre en place, présente des limites : double saisie, risque d'erreurs de retranscription, recherche lente d'un appareil et absence d'indicateurs d'activité.

Dans ce contexte, le stage a eu pour objectif de concevoir et de réaliser \textbf{AirLiquide GMAO+}, une application web de gestion de la maintenance qui remplace ces saisies par un suivi unique, rattaché à une fiche de vie durable pour chaque appareil. Ce rapport décrit la démarche suivie et les résultats obtenus.

Il est organisé en quatre chapitres. Le premier présente l'entreprise et le service d'accueil. Le deuxième décrit le contexte du stage, l'existant et les objectifs. Le troisième détaille les travaux réalisés : analyse, conception et réalisation de l'application. Le quatrième présente les résultats, les difficultés rencontrées, les limites et les perspectives. Une conclusion générale clôt le rapport.

